% appendix_c_flags.tex -- Appendix C Flags and weights. Owner: Appendix writer.
% Sources: Final / Control: slurm/experiments/exp7_final_model.sh header (ENV defaults), :207-257 (settings, RUN,
%   BONE_FIX_ARG, HANDS_ARG, SHAPE_ARG, TIME), :759-771 (the train.py command). Student:
%   slurm/experiments/exp8_distill_vitb.sh header (RECIPE, ENV) and :319-331 (TRAIN_CMD); TIME=20:00:00 from the
%   submission (H 0a, 18 Oct 17:15). Defaults: train.py:1558-1725 and loss.py:440-500 (add_specific_args).
%   Colour-jitter flags: parsed at train.py:1722-1725, read nowhere (grep "args.brightness" over the repo: no hit;
%   data/CLUSTER_ANSWERS.md: no augmentation). Loss weights: tables/T02_loss_weights.tex (not repeated here).
% No job ids and no old run names in this file (plan section 5); Table~\ref{tab:names} maps names to runs.

\chapter{Flags and Weights}
\label{app:flags}

This appendix lists every command-line flag of the training runs of Final and Control (Table~\ref{tab:flags-final})
and of the Student (Table~\ref{tab:flags-student}), with the default in the code, the value used and what the flag
does. The loss weights appear in Table~\ref{tab:loss-weights} of Chapter~\ref{ch:method} and are only named here.
Flags that are not passed keep their default. The commands that start the runs are in Section~\ref{sec:app-commands}.

\begingroup
\expandafter\def\expandafter\UrlBreaks\expandafter{\UrlBreaks\do\_\do\-}%
\footnotesize
\setlength{\tabcolsep}{3pt}
\begin{longtable}{@{}>{\raggedright}p{3.6cm}>{\raggedright}p{1.55cm}>{\raggedright}p{2.4cm}>{\raggedright\arraybackslash}p{4.7cm}@{}}
\caption[Training flags of Final and Control]{Training flags of Final and Control. ``Value'' is Final's; where Control
differs, its value follows in brackets. Paths are relative to the project folder on the cluster. Both runs start from
Base and see 10 epochs of 1,000 iterations at batch 4, that is 40,000 images.}
\label{tab:flags-final}\\
\toprule
Flag & Default & Value (Control) & Meaning \\
\midrule
\endfirsthead
\toprule
Flag & Default & Value (Control) & Meaning \\
\midrule
\endhead
\bottomrule
\endfoot
\multicolumn{4}{@{}l}{\emph{Start and network}} \\
\path{--pretrained} & none & Base checkpoint & weights the run starts from \\
\path{--pretrained_remap} & 1 & 0 & no renaming of weight names; Base was saved by this code \\
\path{--load_only_backbone} & 0 & 0 & load the heads too, not only the backbone \\
\path{--backbone} & ViT-S & \texttt{dinov2\_vitl14} & DINOv2 ViT-L with 14-pixel patches \\
\path{--img_size} & 336 & 672 & input side in pixels (48 $\times$ 48 tokens) \\
\path{--person_center} & head & head & the joint each detection and its predicted position refer to \\
\path{--freeze_backbone} & 0 & 1 & freeze the backbone \ldots \\
\path{--unfreeze_last_n_blocks} & 0 & 4 & \ldots except its last 4 of 24 blocks and final norm \\
\path{--bone_fix} & off & \path{results/experiments/exp6_bone_fix/anny_s2_shape_v5_00099_n200_s0/bone_fix.json} (\texttt{none}) & constant output rotation on neck02, head, metacarpal1.L and finger4-1.R, saved in every checkpoint and re-applied on load (Section~\ref{sec:method-bonefix}) \\
\addlinespace
\multicolumn{4}{@{}l}{\emph{Targets and losses}} \\
\path{--gt_center_joint} & root & \texttt{person\_center} (\texttt{root}) & the joint the depth, translation, detection-cell and offset targets describe: the head, as the network anchors (Control: the pelvis) \\
\path{--alpha_dist}, \path{--alpha_transl} & 1, 1 & 5, 5 (1, 1) & weights of the depth and translation terms (Table~\ref{tab:loss-weights}) \\
\path{--alpha_shape}, \path{--alpha_rotmat}, \path{--alpha_j2d}, \path{--alpha_v2d}, \path{--alpha_global_rot} & 1, 0.1, 1, 1, 0 & 20, 0.1, 0.1, 0.1, 0 & loss weights (Table~\ref{tab:loss-weights}) \\
\path{--boost_neck_weight}, \path{--boost_hand_weight} & 1, 1 & 4, 4 & extra weight on neck, head and hand bones, renormalised to mean 1 (Table~\ref{tab:loss-weights}) \\
\path{--rotmat_loss} & l1 & l1 & rotation term as L1 over the nine matrix entries \\
\path{--global_rot_bones} & head, neck03, wrists & default & bones of the global-rotation term; unused at weight 0 \\
\path{--start_2d_epoch} & 10 & 0 & 2D terms from the first epoch (the epoch counter restarts in every run) \\
\path{--mask_helper_joints_train} & 0 & 1 & hold the 88 helper bones at identity and drop their rotation targets \\
\path{--use_anny_shape} & 0 & 1 & take the phenotype targets from the dataset \\
\path{--mask_3dpw_hands} & 0 & 1 & drop the 38 finger, thumb and metacarpal bones of 3DPW people from the rotation terms \\
\path{--mask_3dpw_shape} & 0 & not passed (0) & shape mask for 3DPW people; on only in Final-masked and Final-partial-mask \\
\path{--debug_anchor} & 0 & 2 & print predicted and target anchors for the first 2 iterations (a check, no effect on training) \\
\addlinespace
\multicolumn{4}{@{}l}{\emph{Data}} \\
\path{--train_data} & BEDLAM & AnnyOne & synthetic training set \\
\path{--train_n}, \path{--val_anny_n} & 0, 0 & 0, 100 & all AnnyOne images except the last 100, which form the holdout \\
\path{--exclude_indices_file} & none & the 55 manifest images & never sample the diagnostic-manifest images; sampling is then with replacement \\
\path{--mix_3dpw_labels} & none & \path{threedpw_anny_v23/labels} & corrected 3DPW-Anny labels mixed into training \\
\path{--mix_3dpw_ratio}, \path{--mix_3dpw_split} & 0.3, train & 0.3, train & 30\,\% of the images of a batch from 3DPW train, on average \\
\path{--val_3dpw_split}, \path{--val_3dpw_n} & validation, 0 & validation, 300 & evaluate on 300 evenly spaced 3DPW validation images at every evaluation \\
\addlinespace
\multicolumn{4}{@{}l}{\emph{Optimisation and bookkeeping}} \\
\path{--batch_size} & 2 & 4 & images per step \\
\path{--n_iters_per_epoch}, \path{--max_iter} & 100, 10,000 & 1,000, 10,000 & 10 epochs of 1,000 steps \\
\path{--learning_rate} & $5\cdot10^{-6}$ & $10^{-5}$ & Adam learning rate \\
\path{--lr_decay_every}, \path{--lr_decay_gamma} & 30, 0.5 & 30, 0.5 & halve every 30 epochs, so constant over the 10 epochs \\
\path{--amp}, \path{--tf32} & 1, 1 & 0, 1 & no mixed precision (it produced NaNs with the Anny layer); TF32 matrix products \\
\path{--num_workers} & 0 & 12 & data-loader processes \\
\path{--eval_freq}, \path{--nb_max_ckpt}, \path{--log_freq} & 1, 10, 100 & 1, 2, 50 & evaluate every epoch, keep the last 2 checkpoints, log every 50 steps \\
\path{--save_dir}, \path{--name} & logs, trainval & run folder, train & output folder \path{results/experiments/exp7_final_model/<run>/train} \\
\addlinespace
\multicolumn{4}{@{}l}{\emph{Job settings of the script (environment variables)}} \\
\path{EVAL_TEST}, \path{N_VIS} & 1, 12 & 1, 12 & after training: holdout trace with 12 panels, then the 3DPW test on the v23 labels \\
\path{TIME}, \path{PARTITIONS} & from the epochs & 5 h; A100, L40S, RTX A6000 & Slurm time limit (60 + 60 + 18 min per epoch) and GPU types \\
\end{longtable}

\begin{longtable}{@{}>{\raggedright}p{3.6cm}>{\raggedright}p{1.55cm}>{\raggedright}p{2.4cm}>{\raggedright\arraybackslash}p{4.7cm}@{}}
\caption[Training flags of the Student]{Training flags of the Student. Flags with the same value as in Final are
grouped in the last block; the targets, masks and loss weights equal Final's so that the Student learns the corrected
depth. 30 epochs of 1,000 iterations at batch 8, that is 240,000 images.}
\label{tab:flags-student}\\
\toprule
Flag & Default & Value & Meaning \\
\midrule
\endfirsthead
\toprule
Flag & Default & Value & Meaning \\
\midrule
\endhead
\bottomrule
\endfoot
\multicolumn{4}{@{}l}{\emph{Student network and its start}} \\
\path{--backbone} & ViT-S & \texttt{dinov2\_vitb14} & DINOv2 ViT-B with 14-pixel patches \\
\path{--pretrained} & none & \path{models/multiHMR/multiHMR_672_B.pt} & NAVER's mesh-recovery ViT-B as the backbone start \\
\path{--pretrained_remap}, \path{--load_only_backbone} & 1, 0 & 1, 1 & rename NAVER's weight names; keep only the backbone \\
\path{--min_backbone_match} & 90 & 90 & stop if less than 90\,\% of the backbone loads \\
\path{--freeze_backbone} & 0 & 0 & train the whole backbone \\
\path{--head_dim} & 0 & 1024 & build the heads at the teacher's width, fed by a 768 to 1024 linear projector \\
\addlinespace
\multicolumn{4}{@{}l}{\emph{Teacher and distillation}} \\
\path{--distill_teacher_ckpt} & none & Final checkpoint & frozen teacher; turns the run into distillation \\
\path{--teacher_bf16} & 0 & 1 & teacher backbone in bfloat16 (no gradients) \\
\path{--init_heads_from_teacher} & 0 & 1 & copy every head tensor of matching name and shape from the teacher \\
\path{--freeze_heads_epochs} & 0 & 5 & keep the copied heads frozen for 5 epochs while backbone and projector adapt \\
\path{--lambda_kd}, \path{--kd_temperature}, \path{--kd_softmax_dim} & 1.0, 4.0, channel & 1.0, 4.0, channel & feature distillation: KL divergence between softened per-token channel distributions (Table~\ref{tab:loss-weights}) \\
\path{--lambda_kd_out} & 0 & 1.0 & output distillation: L1 to the teacher's rotations, phenotypes and depth (Table~\ref{tab:loss-weights}) \\
\path{--bone_fix} & off & the teacher's correction & the four-bone correction of Final during training; after training the Student's own correction is estimated on 200 training images (seed 0) and used for every evaluation \\
\addlinespace
\multicolumn{4}{@{}l}{\emph{Optimisation and bookkeeping}} \\
\path{--batch_size} & 2 & 8 & images per step \\
\path{--n_iters_per_epoch}, \path{--max_iter} & 100, 10,000 & 1,000, 30,000 & 30 epochs of 1,000 steps \\
\path{--learning_rate} & $5\cdot10^{-6}$ & $3\cdot10^{-5}$ & Adam learning rate \\
\path{--lr_decay_every}, \path{--lr_decay_gamma} & 30, 0.5 & 12, 0.5 & halve after epochs 12 and 24 \\
\path{--eval_freq}, \path{--nb_max_ckpt}, \path{--log_freq} & 1, 10, 100 & 5, 3, 50 & evaluate every 5 epochs, keep the last 3 checkpoints \\
\path{--brightness}, \path{--contrast}, \path{--saturation}, \path{--hue} & 0 & 0.2, 0.2, 0.2, 0.05 & parsed but read nowhere in the code: no effect, no augmentation was used \\
\path{TIME} & 14 h & 20 h & Slurm time limit \\
\addlinespace
\multicolumn{4}{@{}l}{\emph{As in Final (Table~\ref{tab:flags-final})}} \\
\multicolumn{4}{@{}p{12.3cm}@{}}{\raggedright \texttt{--train\_data AnnyOne --train\_n 0 --val\_anny\_n 100 --person\_center head --img\_size 672
--mask\_helper\_joints\_train 1 --use\_anny\_shape 1 --alpha\_shape 20 --alpha\_j2d 0.1 --alpha\_v2d 0.1
--boost\_neck\_weight 4 --boost\_hand\_weight 4 --rotmat\_loss l1 --alpha\_rotmat 0.1 --alpha\_global\_rot 0
--start\_2d\_epoch 0 --gt\_center\_joint person\_center --alpha\_dist 5 --alpha\_transl 5 --debug\_anchor 2
--exclude\_indices\_file} (the 55 manifest images) \path{--mix_3dpw_labels} (v23) \texttt{--mix\_3dpw\_ratio 0.3
--mix\_3dpw\_split train --val\_3dpw\_split validation --val\_3dpw\_n 300 --mask\_3dpw\_hands 1 --num\_workers 12
--amp 0 --tf32 1}; no shape mask.} \\
\end{longtable}
\endgroup
